% Public source snapshot of the instructor's Module 11 lecture notes. Executable
% implementations, diagram source, and external images are omitted under the
% boundary in sources/README.md. Headings and claims are preserved for provenance,
% not endorsed as reproduced results; the book re-derives equations and checks
% numerical claims against primary papers before use.
\documentclass{article}
\usepackage{graphicx}
\usepackage{amsmath}
\usepackage{amssymb}
\usepackage{listings}
\usepackage[margin=1in]{geometry}
\usepackage{xcolor}
\usepackage{tikz}
\usetikzlibrary{shapes, arrows, positioning, calc}
\usepackage{tcolorbox}
\usepackage{booktabs}

% Define colors for syntax highlighting
\definecolor{codegreen}{rgb}{0,0.6,0}
\definecolor{codegray}{rgb}{0.5,0.5,0.5}
\definecolor{codepurple}{rgb}{0.58,0,0.82}
\definecolor{backcolour}{rgb}{0.95,0.95,0.92}

% Python style for highlighting
\lstdefinestyle{pythonstyle}{
    backgroundcolor=\color{backcolour},
    commentstyle=\color{codegreen},
    keywordstyle=\color{magenta},
    numberstyle=\tiny\color{codegray},
    stringstyle=\color{codepurple},
    basicstyle=\ttfamily\footnotesize,
    breakatwhitespace=false,
    breaklines=true,
    captionpos=b,
    keepspaces=true,
    numbers=left,
    numbersep=5pt,
    showspaces=false,
    showstringspaces=false,
    showtabs=false,
    tabsize=2
}
\lstset{style=pythonstyle}

\title{Lecture 11.3: Efficient LLM Deployment\\
Quantization, QLoRA, and Running Models at Home}
\author{Deep Learning Class}
\date{\today}

\begin{document}
\maketitle

\begin{tcolorbox}[colback=gray!5!white, colframe=gray!60!black, title=Public Snapshot Boundary]
This provenance snapshot preserves the instructor-authored lecture spine. Executable
implementations, diagram source, and external images have been omitted; their original
locations remain marked in the source. Product and numerical claims are historical
lecture notes and are independently checked before publication in the book.
\end{tcolorbox}

\begin{tcolorbox}[colback=blue!5!white, colframe=blue!75!black, title=Lecture Overview]
\textbf{The democratization challenge}: From billion-parameter models to your laptop

\begin{enumerate}
\item \textbf{The memory bottleneck}: Why 7B models need 66+ GB for training
\item \textbf{Quantization fundamentals}: Representing weights with 4 bits instead of 16
\item \textbf{QLoRA}: Fine-tuning 7B models on consumer GPUs (8--12 GB)
\item \textbf{Deployment strategies}: vLLM, llama.cpp, and MLX for different hardware
\item \textbf{Hands-on}: Fine-tune and deploy Qwen2.5-7B on your machine
\item \textbf{Trade-offs}: Speed, quality, and hardware constraints
\end{enumerate}

\textbf{Practical outcome}: By the end, you will fine-tune and serve a 7B model on consumer hardware!
\end{tcolorbox}

\section{The Memory Problem}

\subsection{The Billion-Parameter Reality}

Modern language models are remarkably capable but extraordinarily large:

\begin{itemize}
\item \textbf{GPT-3}: 175 billion parameters
\item \textbf{LLaMA 2}: 7B, 13B, 70B variants
\item \textbf{Mistral}: 7B parameters (Apache-2.0 license)
\item \textbf{Qwen 2.5}: 7B and 14B variants (Apache-2.0); 72B (Qwen license)
\end{itemize}

\textbf{Question}: Can you run these on your laptop? Let us do the math.

\subsection{Memory Requirements: The Calculation}

A model with $N$ parameters needs memory for:
\begin{enumerate}
\item \textbf{Weights} (the model parameters)
\item \textbf{Activations} (intermediate computations)
\item \textbf{Gradients} (for backpropagation, if training)
\item \textbf{Optimizer states} (momentum, variance for Adam, if training)
\end{enumerate}

\subsubsection{Just Storing Weights}

Each parameter is a number. Standard storage formats:

\begin{center}
\begin{tabular}{lccc}
\toprule
\textbf{Format} & \textbf{Bits/param} & \textbf{Bytes/param} & \textbf{7B model} \\
\midrule
FP32 (float32) & 32 & 4 & 28 GB \\
FP16 (float16) & 16 & 2 & 14 GB \\
BF16 (bfloat16) & 16 & 2 & 14 GB \\
INT8 & 8 & 1 & 7 GB \\
INT4 & 4 & 0.5 & \textbf{3.5 GB} \\
\bottomrule
\end{tabular}
\end{center}

\textbf{Example}: 7 billion parameters in FP16:
\[
7 \times 10^9 \text{ params} \times 2 \text{ bytes/param} = 14 \times 10^9 \text{ bytes} = 14 \text{ GB}
\]

\textbf{Your hardware}:
\begin{itemize}
\item Consumer laptop: 8--16 GB RAM (shared with OS, browser)
\item RTX 3060/4060 GPU: 8--12 GB VRAM
\item MacBook Pro M3: 16--32 GB unified memory
\end{itemize}

\textbf{Conclusion}: A 7B model in FP16 barely fits for \textit{inference}, and definitely not for training!

\subsection{Training Makes It Worse}

During training, memory requirements explode:

\begin{center}
% [Diagram source omitted from the public snapshot; the book uses independently drawn figures.]
\end{center}

\textbf{Why so much}?

\begin{itemize}
\item \textbf{Gradients}: Same size as weights (14 GB for FP16)
\item \textbf{Optimizer states}: Adam stores first and second moments:
\[
\mathbf{m}_t = \beta_1 \mathbf{m}_{t-1} + (1-\beta_1)\mathbf{g}_t, \quad
\mathbf{v}_t = \beta_2 \mathbf{v}_{t-1} + (1-\beta_2)\mathbf{g}_t^2
\]
That is 2 extra copies in FP32 (28 GB!)
\item \textbf{Activations}: Stored during forward pass for backpropagation (10+ GB depending on batch size and sequence length)
\end{itemize}

\textbf{Total for training 7B with Adam in FP16}: $\sim$66 GB

This requires enterprise GPUs (A100 80GB, H100 80GB) costing \$10,000--\$40,000!

\begin{tcolorbox}[colback=red!5!white, colframe=red!50!black, title=The Challenge]
\textbf{Goal}: Enable researchers and practitioners to:
\begin{itemize}
\item \textbf{Run} 7B models for inference on 8 GB GPUs
\item \textbf{Fine-tune} 7B models on 12 GB GPUs
\item \textbf{Deploy} models on CPUs and consumer hardware
\end{itemize}

\textbf{Solution}: Quantization + efficient fine-tuning + optimized serving!
\end{tcolorbox}

\section{Quantization: Fewer Bits, Less Memory}

\subsection{The Core Idea}

\textbf{Quantization}: Represent weights with fewer bits while preserving performance.

\textbf{Analogy}: Image compression
\begin{itemize}
\item Raw photo: 24 bits per pixel (16 million colors)
\item JPEG: Approximate values, much smaller file
\item Visual quality: Often indistinguishable
\end{itemize}

For neural networks:
\begin{itemize}
\item Original weights: FP16 or FP32 (continuous values)
\item Quantized weights: INT8 or INT4 (discrete levels)
\item Model quality: Minimal degradation with proper techniques!
\end{itemize}

\subsection{How Quantization Works}

\textbf{Basic quantization} maps floating-point weights to integers.

\begin{center}
% [Diagram source omitted from the public snapshot; the book uses independently drawn figures.]
\end{center}

\subsubsection{The Quantization Formula}

Given a floating-point weight $w \in \mathbb{R}$, quantize to integer:
\[
w_{\text{quant}} = \text{round}\left(\frac{w - z}{s}\right)
\]

where:
\begin{itemize}
\item $s$ = \textbf{scale} (controls range)
\item $z$ = \textbf{zero-point} (controls offset)
\end{itemize}

To reconstruct (dequantize):
\[
\hat{w} = w_{\text{quant}} \cdot s + z
\]

\subsubsection{Example: 8-bit Quantization}

Weights in range $[-0.5, 0.5]$, quantize to INT8 (256 values):

\begin{enumerate}
\item \textbf{Determine range}: $w_{\min} = -0.5$, $w_{\max} = 0.5$
\item \textbf{Compute scale}:
\[
s = \frac{w_{\max} - w_{\min}}{2^8 - 1} = \frac{1.0}{255} \approx 0.00392
\]
\item \textbf{Quantize} $w = 0.12$:
\[
w_{\text{quant}} = \text{round}\left(\frac{0.12}{0.00392}\right) = \text{round}(30.6) = 31
\]
\item \textbf{Dequantize}:
\[
\hat{w} = 31 \times 0.00392 \approx 0.1215
\]
\end{enumerate}

\textbf{Error}: $|0.12 - 0.1215| = 0.0015$ (very small!)

\subsection{Quantization Strategies}

\begin{center}
\begin{tabular}{lp{5cm}l}
\toprule
\textbf{Method} & \textbf{Description} & \textbf{Bits} \\
\midrule
\textbf{Post-Training (PTQ)} & Quantize after training & 4--8 \\
\textbf{Quantization-Aware (QAT)} & Train with quantization & 4--8 \\
\textbf{Weight-Only} & Quantize weights, keep activations FP16 & 3--4 \\
\textbf{Weight + Activation} & Quantize both & 8 \\
\bottomrule
\end{tabular}
\end{center}

For LLMs, \textbf{weight-only quantization} dominates:
\begin{itemize}
\item Quantize weights to INT4 (16 values)
\item Keep activations in FP16 during computation
\item \textbf{4$\times$ memory reduction} with minimal quality loss!
\end{itemize}

\subsection{NF4: Normal Float 4}

Standard INT4 uses uniform spacing: $\{0, 1, 2, \ldots, 15\}$.

But neural network weights follow a \textbf{normal distribution}!

\textbf{NF4} (used in QLoRA) optimally spaces quantization levels:

\begin{center}
% [Diagram source omitted from the public snapshot; the book uses independently drawn figures.]
\end{center}

\textbf{Key insight}: More quantization levels near zero (high density), fewer in tails (low density).

This gives \textbf{better precision where it matters}!

\subsection{Quantization Methods for LLMs}

\begin{center}
\begin{tabular}{lp{6cm}l}
\toprule
\textbf{Method} & \textbf{Description} & \textbf{Use} \\
\midrule
\textbf{GPTQ} & Post-training using Hessian info & 3--4 bit inference \\
\textbf{AWQ} & Activation-aware (protects salient weights) & Fast inference \\
\textbf{GGUF} & Format for llama.cpp (CPU/Metal) & Laptops, no GPU \\
\textbf{bitsandbytes} & On-the-fly 4/8-bit with NF4 & Training and inference \\
\bottomrule
\end{tabular}
\end{center}

\textbf{Bottom line}: With 4-bit quantization, 7B model: 14 GB $\rightarrow$ \textbf{3.5 GB}!

\section{QLoRA: Fine-Tuning on Consumer GPUs}

\subsection{Recall: LoRA from the Pretraining Lecture}

In the previous lecture, we learned about \textbf{LoRA} (Low-Rank Adaptation):

\begin{itemize}
\item Freeze pre-trained weights $\mathbf{W} \in \mathbb{R}^{d \times d}$
\item Add trainable low-rank adapters: $\Delta \mathbf{W} = \mathbf{B}\mathbf{A}$, where $\mathbf{B} \in \mathbb{R}^{d \times r}$, $\mathbf{A} \in \mathbb{R}^{r \times d}$, and $r \ll d$
\item Forward pass: $\mathbf{h} = \mathbf{W}\mathbf{x} + \mathbf{B}\mathbf{A}\mathbf{x}$
\item Only train $\mathbf{B}$ and $\mathbf{A}$ (tiny fraction of parameters!)
\end{itemize}

\textbf{Memory savings}:
\begin{itemize}
\item Original params: $d^2$
\item LoRA params: $2dr$ (typically $r=8$ or $r=16$, so $2dr \ll d^2$)
\end{itemize}

\textbf{The problem}: Even with LoRA, the base model needs 14 GB in FP16!

\subsection{QLoRA: The Best of Both Worlds}

\textbf{QLoRA} = \textbf{Quantized} base model + \textbf{LoRA} adapters

\begin{center}
% [Diagram source omitted from the public snapshot; the book uses independently drawn figures.]
\end{center}

\textbf{Key innovations in QLoRA}:

\begin{enumerate}
\item \textbf{4-bit NormalFloat (NF4)}: Quantize base to 4 bits optimally
\item \textbf{Double Quantization}: Quantize the quantization constants (extra savings)
\item \textbf{Paged Optimizers}: Handle memory spikes via unified memory
\item \textbf{LoRA adapters in FP16}: Keep trainable params in high precision
\end{enumerate}

\subsection{Memory Breakdown}

For 7B model with QLoRA (rank $r=16$):

\begin{center}
\begin{tabular}{lccc}
\toprule
\textbf{Component} & \textbf{Size} & \textbf{Precision} & \textbf{Trainable?} \\
\midrule
Base weights & 3.5 GB & NF4 (4-bit) & No \\
LoRA adapters & $\sim$50 MB & FP16 & Yes \\
Gradients (adapters) & $\sim$50 MB & FP16 & -- \\
Optimizer states & $\sim$200 MB & FP32 & -- \\
Activations (batch=1) & $\sim$1 GB & FP16 & -- \\
\midrule
\textbf{Total} & \textbf{5--6 GB} & & \\
\bottomrule
\end{tabular}
\end{center}

\textbf{Comparison}:
\begin{itemize}
\item Full fine-tuning (FP16): 66 GB (requires A100 80GB)
\item QLoRA: 5--6 GB (runs on RTX 3060!)
\end{itemize}

\textbf{Performance}: QLoRA achieves 99\% of full fine-tuning quality with \textbf{10$\times$ less memory}!

\subsection{Implementation: Fine-tuning Qwen2.5-7B with QLoRA}

Let us walk through fine-tuning a 7B model on consumer hardware.

\subsubsection{Step 1: Setup and Model Selection}

% [Implementation omitted from the public snapshot: independent provenance was not established. See sources/README.md.]

\textbf{Why Qwen2.5-7B}?
\begin{itemize}
\item Apache-2.0 licensed; redistribution remains subject to its notice and license conditions
\item Strong multilingual capabilities
\item Good long-context performance
\item Active development and community
\end{itemize}

\subsubsection{Step 2: Configure 4-bit Quantization}

% [Implementation omitted from the public snapshot: independent provenance was not established. See sources/README.md.]

\begin{tcolorbox}[colback=yellow!5!white, colframe=orange!50!black, title=What is Double Quantization?]
The quantization parameters (scales $s$) themselves take memory!

\textbf{Double quantization}: Quantize the scales from FP32 to INT8.

For 7B model:
\begin{itemize}
\item Scales in FP32: $\sim$500 MB
\item Scales in INT8: $\sim$125 MB
\item \textbf{Savings}: $\sim$375 MB
\end{itemize}

Small but meaningful when every MB counts!
\end{tcolorbox}

\subsubsection{Step 3: Load Model and Apply LoRA}

% [Implementation omitted from the public snapshot: independent provenance was not established. See sources/README.md.]

\textbf{Observation}: We are training only \textbf{0.55\%} of parameters! Yet this is enough for effective adaptation.

\subsubsection{Step 4: Prepare Dataset}

% [Implementation omitted from the public snapshot: independent provenance was not established. See sources/README.md.]

\subsubsection{Step 5: Train with QLoRA}

% [Implementation omitted from the public snapshot: independent provenance was not established. See sources/README.md.]

\textit{The original implementation and its unverified run-result narration are
omitted from this public snapshot.}

\begin{tcolorbox}[colback=green!5!white, colframe=green!50!black, title=Why QLoRA Works: The Key Insight]
\textbf{Observation}: You do not need high precision everywhere!

\begin{itemize}
\item \textbf{Base weights (frozen)}: Can be 4-bit
\begin{itemize}
\item Fixed during training
\item Dequantized to FP16 for computation
\item 4-bit sufficient for storage
\end{itemize}

\item \textbf{LoRA adapters (trainable)}: Must be FP16
\begin{itemize}
\item Updated by gradients
\item Need precision for training stability
\item But tiny ($<$1\% of parameters)
\end{itemize}
\end{itemize}

\textbf{Result}: 10$\times$ memory reduction, $<$1\% quality loss!
\end{tcolorbox}

\section{Summary and Key Takeaways}

\begin{tcolorbox}[colback=purple!5!white, colframe=purple!50!black, title=What We Learned]
\textbf{Memory bottleneck}:
\begin{itemize}
\item 7B model in FP16: 14 GB (inference), 66+ GB (training)
\item Consumer hardware: 8--16 GB (not enough!)
\end{itemize}

\textbf{Quantization}:
\begin{itemize}
\item FP16 $\rightarrow$ INT4: 4$\times$ memory reduction
\item NF4: Optimal spacing for normal distributions
\item Quality loss: $<$2\% with smart quantization
\end{itemize}

\textbf{QLoRA}:
\begin{itemize}
\item 4-bit base + FP16 LoRA adapters
\item Train 7B models on 8--12 GB GPUs
\item 10$\times$ memory savings, 99\% quality retention
\end{itemize}

\textbf{Deployment}:
\begin{itemize}
\item \textbf{vLLM}: High-throughput GPU serving with PagedAttention
\item \textbf{llama.cpp}: CPU/edge deployment with GGUF format
\item \textbf{MLX-LM}: Optimized for Apple Silicon
\item \textbf{AWQ}: Fast GPU inference with activation-aware quantization
\end{itemize}
\end{tcolorbox}

\subsection{Practical Guidelines}

\textbf{For training}:
\begin{enumerate}
\item Use QLoRA with rank $r=16$ as default
\item Enable gradient checkpointing to save memory
\item Use paged\_adamw\_8bit optimizer
\item Start with learning rate 2e-4 for LoRA
\end{enumerate}

\textbf{For deployment}:
\begin{enumerate}
\item Profile your use case (latency vs throughput)
\item Choose quantization level (Q4\_K\_M is good default)
\item Benchmark on your specific hardware
\item Monitor quality degradation on held-out test set
\end{enumerate}

\subsection{Common Pitfalls}

\begin{itemize}
\item \textbf{Over-quantization}: 3-bit often too aggressive, stick to 4-bit
\item \textbf{Wrong target modules}: Always include Q, K, V, O projections in LoRA
\item \textbf{Insufficient calibration data}: AWQ needs representative samples
\item \textbf{Memory leaks}: Clear KV-cache between generation batches
\item \textbf{Ignoring gradient checkpointing}: Can save 30--50\% memory during training
\end{itemize}

\section{Exercises and Further Exploration}

\subsection{Exercise 1: Quantization Trade-offs}

Quantize Qwen2.5-7B to different precisions and evaluate:

\begin{enumerate}
\item Quantize to: FP16, INT8, INT4 (NF4), INT4 (uniform), INT3
\item Evaluate perplexity on held-out test set
\item Measure inference speed (tokens/second)
\item Plot: bits per parameter vs. perplexity
\end{enumerate}

\textbf{Question}: Where is the sweet spot for your application?

\textbf{Deliverable}: Report with plots and recommendations.

\subsection{Exercise 2: LoRA Rank Ablation}

Fine-tune with different ranks: $r \in \{4, 8, 16, 32, 64\}$

Measure:
\begin{itemize}
\item Training time per epoch
\item Peak memory usage
\item Validation loss
\item Task-specific metrics (accuracy, F1, etc.)
\end{itemize}

\textbf{Question}: What is the optimal rank for your task? Does it depend on model size?

\subsection{Exercise 3: Deployment Comparison}

Deploy your fine-tuned model with multiple methods:

\begin{center}
\begin{tabular}{lcccc}
\toprule
\textbf{Method} & \textbf{Setup time} & \textbf{Latency} & \textbf{Throughput} & \textbf{Memory} \\
\midrule
HF + bitsandbytes & ? & ? & ? & ? \\
vLLM & ? & ? & ? & ? \\
llama.cpp (CPU) & ? & ? & ? & ? \\
llama.cpp (Metal) & ? & ? & ? & ? \\
MLX-LM & ? & ? & ? & ? \\
\bottomrule
\end{tabular}
\end{center}

\textbf{Test conditions}:
\begin{itemize}
\item Single request latency (time to first token, total time)
\item Batch throughput (tokens/second with batch sizes 1, 8, 32)
\item Memory usage (idle, during inference)
\end{itemize}

\textbf{Deliverable}: Comprehensive benchmark report with recommendations for different use cases.

\subsection{Further Reading}

\textbf{Core papers}:
\begin{itemize}
\item Dettmers et al., ``QLoRA: Efficient Finetuning of Quantized LLMs'' (2023)
\item Frantar et al., ``GPTQ: Accurate Post-Training Quantization for GPT Models'' (2022)
\item Lin et al., ``AWQ: Activation-aware Weight Quantization for LLM Compression'' (2023)
\item Kwon et al., ``Efficient Memory Management for Large Language Model Serving with PagedAttention'' (2023)
\item Dettmers et al., ``LLM.int8(): 8-bit Matrix Multiplication for Transformers at Scale'' (2022)
\end{itemize}

\textbf{Tools and frameworks}:
\begin{itemize}
\item vLLM documentation: \texttt{docs.vllm.ai}
\item llama.cpp: \texttt{github.com/ggerganov/llama.cpp}
\item AutoAWQ: \texttt{github.com/casper-hansen/AutoAWQ}
\item MLX: \texttt{ml-explore.github.io/mlx}
\item bitsandbytes: \texttt{github.com/TimDettmers/bitsandbytes}
\end{itemize}

\section{Looking Ahead}

\subsection{Emerging Techniques}

The field is rapidly evolving. Recent developments include:

\begin{itemize}
\item \textbf{GPTQ-R}: Refined GPTQ with better Hessian approximation
\item \textbf{QuIP}: Quantization with Incoherence Processing (2-bit!)
\item \textbf{SpQR}: Sparse Quantized Representation
\item \textbf{SmoothQuant}: Migrate difficulty from activations to weights
\item \textbf{EETQ}: Efficient 8-bit quantization for training
\end{itemize}

\subsection{Hardware Trends}

New hardware is making efficient inference easier:

\begin{itemize}
\item \textbf{NPUs}: Neural Processing Units in consumer devices
\item \textbf{Qualcomm AI Engine}: On-device LLM inference on phones
\item \textbf{Apple Neural Engine}: M-series optimizations
\item \textbf{AMD ROCm}: Expanding GPU ecosystem beyond NVIDIA
\end{itemize}

\subsection{The Democratization Vision}

\begin{tcolorbox}[colback=blue!5!white, colframe=blue!50!black, title=The Future is Accessible]
\textbf{Three years ago}: Running LLMs required million-dollar infrastructure.

\textbf{Today}: With the techniques you learned:
\begin{itemize}
\item Run 7B models on \$300 consumer GPUs
\item Fine-tune on laptops with 16 GB RAM
\item Deploy on edge devices and phones
\end{itemize}

\textbf{Tomorrow}: Even more efficient methods will emerge!
\begin{itemize}
\item Better quantization (sub-4-bit without quality loss)
\item Sparse models (Mixture-of-Experts at edge)
\item Neural architecture search for efficiency
\item Hardware-software co-design
\end{itemize}

\textbf{Your opportunity}: Use these tools to build LLM applications that were impossible just a few years ago!
\end{tcolorbox}

\section{Final Thoughts}

The techniques you learned today are not just academic exercises---they are \textbf{actively used in production}:

\begin{itemize}
\item \textbf{Startups}: Deploy LLM services without massive GPU clusters
\item \textbf{Research labs}: Enable more scientists to experiment with large models
\item \textbf{Edge AI}: Run models on robots, drones, IoT devices
\item \textbf{Privacy applications}: Keep sensitive data local with on-device inference
\item \textbf{Developing regions}: Access AI without expensive cloud services
\end{itemize}

As models grow even larger (GPT-4 rumored at 1.7T parameters, future models may reach 10T+), efficiency techniques become \textit{more} critical, not less. The race is not just about bigger models---it is about making powerful models accessible, deployable, and sustainable.

\vspace{1cm}

\begin{center}
\textit{``The best model is not the largest model---}\\
\textit{it is the model you can actually deploy and use.''}
\end{center}

\vspace{1cm}

\hrule

\vspace{0.5cm}

\textbf{Next lecture}: Generative Models and Future Frontiers (VAEs, GANs, Diffusion Models)

\textbf{Lab assignment}:
\begin{enumerate}
\item Fine-tune Qwen2.5-7B on instruction-following task using QLoRA
\item Convert to GGUF and test on CPU
\item Deploy with vLLM and benchmark performance
\item Submit: Code, benchmarking report, and deployment guide
\end{enumerate}

\textbf{Due}: Two weeks from today

\end{document}
